\documentclass[11pt, letterpaper]{article}

\usepackage[margin=1in]{geometry}
\usepackage{amsmath}
\usepackage{booktabs}
\usepackage{array}
\usepackage{enumitem}
\usepackage{parskip}
\usepackage{xcolor}
\usepackage{mdframed}
\usepackage{hyperref}
\usepackage{titlesec}
\usepackage{multicol}

% Colors
\definecolor{blue}{HTML}{1e40af}
\definecolor{lightblue}{HTML}{dbeafe}
\definecolor{boxblue}{HTML}{eff6ff}
\definecolor{ruleblue}{HTML}{3b82f6}
\definecolor{yellow}{HTML}{fef9c3}
\definecolor{ruleyellow}{HTML}{ca8a04}
\definecolor{textgray}{HTML}{334155}
\definecolor{mutedgray}{HTML}{64748b}

% Section formatting
\titleformat{\section}{\large\bfseries\color{blue}}{}{0em}{}[\titlerule]
\titleformat{\subsection}{\normalsize\bfseries\color{blue}}{}{0em}{}
\titlespacing{\section}{0pt}{18pt}{8pt}
\titlespacing{\subsection}{0pt}{12pt}{4pt}

% Blue info box
\newmdenv[
  backgroundcolor=boxblue,
  linecolor=ruleblue,
  linewidth=3pt,
  topline=false,
  rightline=false,
  bottomline=false,
  innerleftmargin=10pt,
  innerrightmargin=10pt,
  innertopmargin=8pt,
  innerbottommargin=8pt,
  skipabove=10pt,
  skipbelow=10pt,
]{infobox}

% Yellow warning box
\newmdenv[
  backgroundcolor=yellow,
  linecolor=ruleyellow,
  linewidth=3pt,
  topline=false,
  rightline=false,
  bottomline=false,
  innerleftmargin=10pt,
  innerrightmargin=10pt,
  innertopmargin=8pt,
  innerbottommargin=8pt,
  skipabove=10pt,
  skipbelow=10pt,
]{warnbox}

\hypersetup{colorlinks=true, urlcolor=blue, linkcolor=blue}

% Tighter list spacing
\setlist[itemize]{noitemsep, topsep=4pt}
\setlist[enumerate]{noitemsep, topsep=4pt}

\begin{document}

% Header
{\LARGE\bfseries\color{blue} Groundwater Commons Game}\\[4pt]
{\color{mutedgray} Student Handout \quad|\quad Environmental Economics}\\[-4pt]
\noindent\rule{\linewidth}{1pt}

\vspace{6pt}

%------------------------------------------------------------------
\section{Overview}

You are a farmer who pumps groundwater to irrigate your crops. The water comes from a shared underground aquifer. You will make pumping decisions across several rounds and try to maximize your total profit --- but so will everyone else sharing your aquifer. Your choices affect the water level that \emph{everyone} faces in future rounds.

\begin{infobox}
\textbf{Big idea:} This game illustrates the \textbf{tragedy of the commons} --- a situation where individually rational decisions lead to a collectively worse outcome for a shared resource.
\end{infobox}

%------------------------------------------------------------------
\section{How the Aquifer Works}

The aquifer has a \textbf{stock level} $S$ that changes every round based on how much water everyone pumps and how much naturally recharges:

\begin{infobox}
\[
S_{t+1} = \max\!\left\{0,\; S_t + R - \textstyle\sum_{i=1}^{n} q_{i,t}\right\}
\]
\vspace{-6pt}
\begin{itemize}
  \item $S_t$ = current water stock in the aquifer
  \item $R$ = natural recharge per round (constant)
  \item $q_{i,t}$ = pumping by player $i$ in round $t$; \;$n$ = number of players
\end{itemize}
\end{infobox}

If the group pumps more than the aquifer recharges, the stock declines. If the stock reaches zero, pumping becomes very expensive and profits collapse. The game ends after 8 rounds (or when the aquifer is exhausted).

%------------------------------------------------------------------
\section{How Your Profit Is Calculated}

Each round, your profit depends on how much you pump ($q$) and the current water level ($S$):

\begin{infobox}
\[
\pi(q, S) \;=\; \underbrace{P\,q}_{\text{revenue}}
           \;-\; \underbrace{\tfrac{\gamma}{2}\,q^2}_{\text{congestion}}
           \;-\; \underbrace{c\,(S_{\max}-S)\,q}_{\text{depth cost}}
\]
\end{infobox}

\vspace{6pt}

\begin{tabular}{@{}p{4.5cm} p{6.5cm} p{2.5cm}@{}}
\toprule
\textbf{Term} & \textbf{What it represents} & \textbf{Default} \\
\midrule
$P \cdot q$ & Revenue from selling your water & $P = 10$ \\[4pt]
$\dfrac{\gamma}{2} \cdot q^2$ & Congestion cost --- pumping too much at once has diminishing returns & $\gamma = 0.08$ \\[8pt]
$c (S_{\max}-S) q$ & Depth cost --- the lower the water table, the more expensive to pump & $c = 0.006$ \\
\bottomrule
\end{tabular}

\vspace{6pt}

\begin{warnbox}
\textbf{Key insight:} When the aquifer depletes, pumping costs rise for \emph{everyone}. This is the \textbf{stock externality} --- your pumping today imposes costs on all players (including yourself) in future rounds.
\end{warnbox}

%------------------------------------------------------------------
\section{Game Modes}

\begin{multicols}{2}
\subsection*{Solo Mode}
You control all 6 wells yourself. Set the same pumping rate for each well and observe how the aquifer responds over 8 rounds. Use this to explore the dynamics \emph{before} playing with your group.

\textit{Try pumping at full capacity every round, then try pumping conservatively. What do you notice?}

\columnbreak

\subsection*{Group Mode}
Each student controls one well in a shared aquifer. Every round you independently choose your pumping rate, and the round advances once everyone submits. The player with the highest \textbf{cumulative profit} at the end wins.
\end{multicols}

%------------------------------------------------------------------
\section{How to Play --- Group Mode}

\begin{enumerate}
  \item Navigate to the game website and click \textbf{Group Mode}.
  \item One person acts as \textbf{host}: click ``Host Controls,'' set the number of players, then click \textbf{Create Room}. Share the 5-character room code with your group.
  \item All other players click ``Join as Player,'' enter the room code and your name, then click \textbf{Join Room}.
  \item Once everyone has joined, the host clicks \textbf{Start Game}.
  \item Each round: move the slider to set your pumping rate ($0$ to $q_{\max}$), then click \textbf{Finalize Pumping Decision}. The round advances automatically when all players have submitted.
  \item After 8 rounds, a final leaderboard shows cumulative profits for all players.
\end{enumerate}

%------------------------------------------------------------------
\section{Discussion Questions}

\begin{itemize}
  \item What pumping rate maximizes your profit in a \emph{single} round, ignoring future costs?
  \item What pumping rate would be best for the \emph{group as a whole} over all 8 rounds?
  \item Why might these two answers differ? What happens to the aquifer if everyone pumps to maximize their individual round profit?
  \item Did your group implicitly cooperate, or did you experience overextraction? How did the aquifer stock change over time?
  \item What real-world policies or institutions could help groups avoid the tragedy of the commons?
\end{itemize}

%------------------------------------------------------------------
\section{Key Terms}

\begin{tabular}{@{}p{3.8cm} p{10cm}@{}}
\toprule
\textbf{Term} & \textbf{Definition} \\
\midrule
Stock externality & Each unit you extract today lowers the aquifer level, raising pumping costs for everyone in every future round. \\[4pt]
Congestion externality & Pumping a large quantity within a single round faces increasing per-unit costs (captured by the $\gamma$ term). \\[4pt]
Nash equilibrium & The outcome where each player is choosing their best response to everyone else's choices --- often leads to overextraction here. \\[4pt]
Social optimum & The pumping plan that maximizes total group profit across all rounds, accounting for future stock effects. \\[4pt]
Tragedy of the commons & When individually rational behavior leads to collective overuse and degradation of a shared resource. \\
\bottomrule
\end{tabular}

\vspace{20pt}
\noindent\rule{\linewidth}{0.4pt}\\[3pt]
{\small\color{mutedgray} Groundwater Commons Game \quad\textbullet\quad \textcopyright\ 2025 Jenna Anders. All rights reserved.}

\end{document}
